\chapter{Dieu et le Christ dans les épitaphes}\label{ch:dieu_christ}

\section{Inventaire}

Dieu est désigné, sans restitution, dans huit textes de l’ensemble C, et peut-être dans un neuvième : \emph{deus} au \nnum{73} (\emph{deo dica⟨vi⟩sti}) et au \nnum{475} (\emph{dedit deus deus tulit}), \emph{ancilla dei} aux \nnums{214}, 306 et 458, \emph{homo dei} au \nnum{474}, \emph{dominus} aux \nnums{85} (\emph{omnia domino}) et 400 (\emph{in luce domini}), et peut-être au \nnum{232} (\emph{in dom(ino?)}). Le Christ est nommé en toutes lettres dans un seul texte (\emph{vitam Cristi}, \nnum{306}) ; il l’est sous une forme abrégée dans quatre autres (\emph{in pace Chr(isti)}, \nnums{140} et 151 ; \emph{in Chr(isto)}, \nnum{311} ; \emph{ancilla dei et Chr(ist)i}, \nnum{458}). Le nom de Jésus n’apparaît dans aucun texte.

Ces désignations sont donc à la fois rares et peu variées. Elles se répartissent en trois usages : la qualification du défunt (\emph{ancilla dei}, \emph{homo dei}), la formule du repos (\emph{in pace Christi}, \emph{in Christo}), l’éloge et la consolation.

\section{Le monogramme et le mot}

La formule \emph{in pace Christi} pose un problème de transcription que la confrontation des témoins rend visible. Au \nnum{38}, l’épitaphe de Privata (352), l’EDCS transcrit \emph{in pace Chr(isti)}, l’EDB \emph{in pace ((monogramma Christi))}, et l’estampage conservé dans la collection de l’Ohio State University sous le numéro ILCV 3252A montre un christogramme après IN PACE. Deux témoins indépendants de l’EDCS concordent : la pierre porte un monogramme, non le mot écrit en lettres.

Aux \nnums{140} et 151 (\emph{in pace Chr(isti)}) et au \nnum{311} (\emph{in Chr(isto)}), les deux bases transcrivent au contraire une abréviation, mais l’EDB décrit un chrisme parmi les données matérielles. Il est probable que, dans ces trois textes, le développement éditorial rende le même monogramme employé comme compendium du nom ; il n’est pas exclu qu’il s’agisse d’une abréviation en lettres accompagnée d’un chrisme décoratif. Seules les photographies permettraient de trancher. Au \nnum{458}, en revanche, l’EDB transcrit \emph{Chr(ist)i} : la lettre finale I est gravée, ce qui désigne une abréviation en lettres.

La différence n’est pas seulement philologique. Un monogramme employé comme abréviation fait partie de la phrase : \emph{in pace Christi} se lit, que le nom soit écrit en lettres ou en signe. Un monogramme juxtaposé n’appartient pas à la phrase : il marque le monument. Les profils du chapitre~\ref{ch:typologie} en tiennent compte en écartant ces quatre textes du profil strict.

\section{\emph{Ancilla dei}, \emph{homo dei} : qualifier la personne}

Trois textes qualifient une défunte d’\emph{ancilla dei} (\nnums{214}, 306 et 458 ; ce dernier avec \emph{ancilla dei et Chr(isti)}), un quatrième un défunt d’\emph{homo dei} (\nnum{474}). Ces expressions désignent l’appartenance de la personne à Dieu. Dans deux cas au moins, elles sont associées à la virginité : Maximilla, \emph{virgo ancilla dei} (\nnum{306}, ligne dont la présence doit être contrôlée sur l’édition, car elle manque dans le texte de l’EDB), et Praetiosa (\emph{Praetioso} dans l’EDCS), \emph{puella virgo XII tantum ancilla dei et Chr(isti)}, morte à douze ans (\nnum{458}). Au \nnum{474}, \emph{homo dei} ouvre une série de participes ascétiques (\emph{confirmans trinitatem}, \emph{amans castitatem}, \emph{respuens mundum}). Au \nnum{214}, l’\emph{ancilla dei} est une femme qui « a possédé cette maison de ses propres biens » et à qui l’on demande de prier pour l’enfant qu’elle laisse.

La formule \emph{ancilla dei et Chr(isti)} du \nnum{458}, que le catalogue de travail classait comme « binaire (Père et Fils) », associe Dieu et le Christ dans une formule de dévotion. Elle ne contient ni le mot « Père » ni le mot « Fils » et ne peut être tenue, sans argument supplémentaire, pour l’expression d’une doctrine des relations divines.

\section{Dieu dans l’éloge et dans l’épreuve}

Trois textes intègrent Dieu à un éloge ou à une consolation.

Le poème du \nnum{73} (363) loue Eleuteria, morte à environ 75 ans, qui « s’est toujours consacrée à Dieu » (\emph{quae te semper deo dica⟨vi⟩sti}, lu). La crainte de Dieu (\emph{[quaeque deu]m metuisti}) n’est qu’une restitution de l’EDCS, que l’EDB ne partage pas (\emph{[quaeque de]metuisti}).

Au \nnum{400} (397), l’esprit de Severianus « a été accueilli dans la lumière du Seigneur » (\emph{cuius spiritus in luce domini susceptus est}). La phrase est lue dans les deux bases, et l’EDB décrit sur la plaque un chrisme et deux croix monogrammatiques, dont l’une est placée immédiatement après \emph{susceptus est}.

Le \nnum{475} (403) reprend, sous une forme libre, la parole de Job : \emph{dedit deus, deus tulit}. La Vulgate porte \emph{Dominus dedit, Dominus abstulit} (Jb 1, 21) ; la forme épigraphique, avec \emph{deus} et \emph{tulit}, peut dépendre d’une version latine antérieure ou d’une citation de mémoire, ce que ce seul texte ne permet pas de décider.

\section{Le Christ dans les textes poétiques}

Deux textes poétiques nomment le Christ au-delà du formulaire. L’épitaphe de Maximilla (\nnum{306}, 389) commence par l’éloge d’une femme « capable de garder l’alliance du Christ » (\emph{poten[s servare foedera C]hristi}, restitution partielle), et poursuit : \emph{sed melius vivit vitam Cristi qui sperat aeternam}, « mais il vit mieux, celui qui espère la vie éternelle du Christ » (lu dans les deux bases). L’épitaphe d’Eutychius (\nnum{352}, 393) dit qu’« il possède la récompense de la lumière, croyant au Christ » (\emph{[qui Ch]ristum credens praemia lucis habet}) ; le début du nom est restitué, mais la fin \emph{-ristum} et le participe \emph{credens} sont lus.

\section{Ce que montrent ces textes}

Dans les épitaphes datées de 340–403, Dieu et le Christ sont nommés pour qualifier le défunt, pour situer son repos et pour consoler. Aucun texte ne développe un énoncé sur leur nature ou leur relation. Le Christ y est l’objet de la foi et le lieu du repos (\emph{in Christo}, \emph{in pace Christi}), le garant de la vie éternelle (\nnum{306}) et de la lumière promise (\nnum{352}). Cette pauvreté de l’énoncé ne permet aucune inférence sur la christologie des commanditaires : elle est d’abord un fait du genre épigraphique (chapitre~\ref{ch:rarete}).
